\section{Conclusion}

This report combined analytical calculations, process and device simulations, cleanroom processing, and electrical measurements to study the BICMOS5 process and the behaviour of the fabricated NMOS and PMOS devices. Overall, the results show good qualitative agreement between theory, Sentaurus simulations, and the experimental data, while the remaining differences can be explained by simplifying assumptions in the analytical models, the full thermal budget of the real process, and measurement-structure-dependent effects. 

From the technology assignments and process simulations, the main trends of implantation, diffusion, and masking were confirmed. In crystalline silicon, channeling strongly affects implantation depth and profile tails, whereas in amorphous silicon the dopant profiles become much closer to Gaussian. The simulations also showed that arsenic is the most suitable dopant for shallow, highly doped NMOS source/drain regions, while phosphorus is more suitable for the N-well because it diffuses deeper into the substrate. In addition, the masking simulations confirmed that a thin photoresist or oxide layer is insufficient for the chosen phosphorus implant, while nitride provides the highest stopping power, although it is less practical as a standard implantation mask.

Comparison of analytical junction-depth calculations with Sentaurus process simulation showed clear differences. The simulated junction depths were approximately $0.537\,\mu\text{m}$ for SN, $2.085\,\mu\text{m}$ for the N-well, and $0.594\,\mu\text{m}$ for SP, whereas the simple Chapter 2 calculations predicted much shallower SN and SP junctions and a slightly deeper N-well. This difference is expected because the analytical treatment assumes simplified Gaussian or limited-source profiles and neglects the later thermal steps of the full CMOS process, while the simulation includes dopant redistribution during oxidation and annealing. 

The processing assignments confirmed the expected behaviour of oxide growth, etching, and metal patterning. Wet oxidation produced a more uniform oxide than PECVD, while PECVD oxide etched significantly faster in buffered HF because of its lower density and higher porosity. For the aluminium etch, wet etching showed high selectivity to the underlying oxide but caused isotropic underetching, whereas dry etching gave straighter and better-defined metal lines at the cost of lower oxide selectivity. These observations are consistent with the expected trade-off between selectivity and dimensional control in wet and dry processing. 

The electrical measurements further validated the simulation results. Van der Pauw measurements gave sheet resistances of about $1740\,\Omega/\square$ for the N-well, $60.2\,\Omega/\square$ for SN, and $488\,\Omega/\square$ for SP in N-well, and the extracted average active dopant concentrations were in good agreement with the earlier Chapter 4 values. The ELM measurements showed clear lateral out-diffusion, which was small for SN, about $0.15\,\mu\text{m}$ per side, but much larger for the N-well, on the order of $0.8$ to $2.1\,\mu\text{m}$ per side, consistent with the longer and hotter phosphorus drive-in. This confirms that lateral diffusion becomes increasingly important for lightly doped and deeply driven layers.

The transistor measurements also matched the expected device physics. For long-channel devices, the extracted threshold voltages and mobilities remained nearly constant with geometry, giving electron mobilities around $750$--$800\,\text{cm}^2/\text{Vs}$ and hole mobilities around $230$--$300\,\text{cm}^2/\text{Vs}$. Short-channel devices showed threshold-voltage roll-off, higher apparent mobility, stronger channel-length modulation, and in the PMOS case even the onset of punch-through, which agrees with the trends predicted from device theory and seen in the simulations.

Finally, the influence of the VT-adjust implant was clearly demonstrated. Increasing the boron VT-adjust dose increased the NMOS threshold voltage monotonically from about $-0.034\,\text{V}$ to $0.819\,\text{V}$, while the PMOS threshold voltage became less negative, changing from about $-3.702\,\text{V}$ to $-3.300\,\text{V}$. This confirms the theoretical expectation that the boron implant raises the channel acceptor concentration for NMOS and partially compensates the N-well surface doping for PMOS. Therefore, the measurements, simulations, and theoretical analysis consistently show how implantation and thermal processing determine the final electrical behaviour of the BICMOS5 devices.